%% From the Extended Hamming Code to the E8 Lattice in Lean
%% LaTeX version generated from Hamming_ConstructionA_E8_Manuscript_Revision.md

\documentclass[12pt]{article}

\usepackage[margin=1.1in]{geometry}
\usepackage{amsmath,amssymb,amsthm}
\usepackage{mathtools}
\usepackage{booktabs}
\usepackage{array}
\usepackage{longtable}
\usepackage{listings}
\usepackage{xcolor}
\usepackage{hyperref}
\usepackage{microtype}
\usepackage{enumitem}

%% -----------------------------------------------------------------------
%% Lean / code listing style
%% -----------------------------------------------------------------------
\definecolor{leanblue}{rgb}{0.15,0.15,0.6}
\definecolor{leangray}{rgb}{0.4,0.4,0.4}
\definecolor{leangreen}{rgb}{0.1,0.45,0.1}

\lstdefinestyle{lean}{
  basicstyle=\small\ttfamily,
  keywordstyle=\color{leanblue}\bfseries,
  commentstyle=\color{leangray}\itshape,
  stringstyle=\color{leangreen},
  breaklines=true,
  breakatwhitespace=true,
  columns=fullflexible,
  keepspaces=true,
  frame=none,
  xleftmargin=1.5em,
}

\lstdefinestyle{shell}{
  basicstyle=\small\ttfamily,
  breaklines=true,
  frame=none,
  xleftmargin=1.5em,
}

%% -----------------------------------------------------------------------
%% Theorem environments
%% -----------------------------------------------------------------------
\newtheorem{theorem}{Theorem}[section]
\newtheorem{proposition}[theorem]{Proposition}
\newtheorem{lemma}[theorem]{Lemma}
\newtheorem{corollary}[theorem]{Corollary}
\theoremstyle{definition}
\newtheorem{definition}[theorem]{Definition}
\newtheorem{remark}[theorem]{Remark}

%% -----------------------------------------------------------------------
%% Macros
%% -----------------------------------------------------------------------
\newcommand{\Z}{\mathbb{Z}}
\newcommand{\Q}{\mathbb{Q}}
\newcommand{\R}{\mathbb{R}}
\newcommand{\C}{\mathbb{C}}
\newcommand{\F}{\mathbb{F}}
\newcommand{\N}{\mathbb{N}}
\newcommand{\code}[1]{\texttt{#1}}
\newcommand{\lean}[1]{\texttt{#1}}
\newcommand{\sqnorm}[1]{\|#1\|^2_{\mathrm{sq}}}
\newcommand{\hammingwt}{\mathrm{wt}}
\newcommand{\Theta}{\vartheta}
\newcommand{\E}[1]{E_{#1}}

\hypersetup{
  colorlinks=true,
  linkcolor=leanblue,
  urlcolor=leanblue,
  citecolor=leanblue,
}

%% -----------------------------------------------------------------------
\title{From the Extended Hamming Code to the E8 Lattice in Lean}
\author{}
\date{}

\begin{document}
\maketitle

%% -----------------------------------------------------------------------
\begin{abstract}
We formalize in Lean~4 the classical construction of the E8 lattice from the
extended binary Hamming code.  Starting from an explicit parity-check matrix
for the $[8,4,4]$ code, we prove that the code is Type~II and unique among
binary linear $[8,4,4]$ codes up to coordinate permutation.  We then apply
Construction~A, exhibit an explicit integer basis, compute its Gram form, and
prove that the scaled lattice is even and unimodular.  The same concrete model
also gives the standard E8 metric data: minimum squared norm~$2$, exactly $240$
minimal vectors, an explicit bridge to an independently formalized E8 root list,
and a Cartan-matrix comparison.

The formalization also connects the coding-theoretic model to the theta-series
description of E8.  In the standalone \code{CodeLatticeE8} package, it proves
the semantic $q^0$ and $q^1$ theta coefficients and a finite Hamming
weight-contribution calculation agreeing with the normalized Eisenstein
coefficients through $q^6$.  The all-coefficient analytic route is isolated in
a Sphere-Packing-Lean-facing package plan: it transports
Construction~A shell counts to the rank-eight theta modular form and compares
that modular form with the normalized Eisenstein series $E_4$, giving the
classical identity
\[
  \Theta_{E_8}(q) = E_4(q) = 1 + 240\sum_{n>0}\sigma_3(n)\,q^n
\]
when the bridge is built.  The development also compares the local
Construction~A basis with the E8 basis used by Sphere-Packing-Lean, the Lean
formalization of Viazovska's dimension-eight sphere-packing theorem.  The
result is a verified bridge from a finite coding-theoretic object to the E8
lattice appearing at the analytic endpoint of the formalized sphere-packing
proof.
\end{abstract}

%% -----------------------------------------------------------------------
\section{Introduction}
\label{sec:intro}

The E8 lattice admits many classical descriptions.  It can be introduced as
the unique even unimodular positive-definite lattice of rank eight~\cite{ConwaySloane1999},
as the root lattice of the exceptional root system $E_8$~\cite{Bourbaki2002},
as a half-integer lattice in $\R^8$~\cite{ConwaySloane1999}, as the lattice of
Coxeter integral octonions, or as the
Construction~A lattice attached to the extended binary Hamming
code~\cite{LeechSloane1971,ConwaySloane1999}.  These descriptions are equivalent
in the mathematical literature, but a formal proof assistant asks for explicit
maps, explicit normalizations, and kernel-checked statements connecting them.

This paper formalizes the coding-theoretic route:
\begin{center}
\begin{tabular}{l}
extended binary Hamming code $[8,4,4]$ \\
$\quad\to$ Construction~A integer lattice \\
$\quad\to$ scaled even unimodular rank-eight lattice \\
$\quad\to$ $240$ minimal vectors \\
$\quad\to$ E8 root and Cartan data \\
$\quad\to$ finite theta coefficients,\\
$\quad\quad$ with Sphere-Packing-Lean route to $\Theta_{E_8} = E_4$ \\
$\quad\to$ Sphere-Packing-Lean E8 basis comparison
\end{tabular}
\end{center}

The contribution is complementary to Sphere-Packing-Lean~\cite{SPL2025}.
Sphere-Packing-Lean formalizes the analytic proof that the E8 sphere packing
is optimal in dimension eight~\cite{Viazovska2017}.  The present development
supplies a finite and combinatorial entrance to the same E8 object: it starts
from the extended Hamming code, builds the lattice by Construction~A, and
identifies the resulting object with the standard E8 data used elsewhere in
the Lean ecosystem.

The main Lean-facing theorem index for the local Hamming-to-E8 development is
\code{CodeLatticeE8.Publication.TheoremIndex}.  The theorem spine is
organized through the standalone package root
\code{CodeLatticeE8} and the optional Sphere-Packing-Lean (SPL) root \code{CodeLatticeE8SPL}.

\subsection{Contributions}

The formalization proves the following.
\begin{enumerate}
\item The explicitly defined extended binary Hamming code is Type~II: it is
      self-dual and doubly even.

\item Every binary linear $[8,4,4]$ code is equivalent to the extended Hamming
      code under a coordinate permutation.

\item The Construction~A lattice has an explicit integer basis, and its scaled
      Gram determinant is $1$.

\item Every lattice vector has the parity and norm divisibility expected of the
      E8 normalization; in particular, the scaled lattice is even.

\item The scaled minimum squared norm is $2$.

\item The list of minimal vectors has exactly $240$ elements and is complete.

\item These $240$ minimal vectors map to the repository's independently
      formalized E8 root list.

\item An explicit unimodular change of basis relates the Construction~A Gram
      matrix to the E8 Cartan matrix~\cite{Bourbaki2002,ConwaySloane1999}.

\item The standalone package proves the semantic zero and first theta
      coefficients, a finite Hamming weight-contribution table check through
      $q^6$, and an all-shell semantic Construction~A convolution theorem.
      The all-coefficient analytic identity is completed in the
      SPL-facing root.

\item The local basis model is compared with the E8 model used by
      Sphere-Packing-Lean, giving a formal bridge from the code-built lattice
      to the analytic sphere-packing endpoint.
\end{enumerate}

\subsection{Scope}

We do not reprove Viazovska's analytic sphere-packing theorem.  We also do not
use the abstract classification theorem saying that there is a unique even
unimodular positive-definite lattice of rank eight.  Instead, the development
works concretely: basis matrices, Gram matrices, explicit shell counts, explicit
roots, and explicit transition maps carry the construction from the code to E8.

This concrete approach has two advantages.  First, it gives a small,
reviewable formal proof of the code-to-lattice route.  Second, it makes all
normalization choices visible, which is essential when comparing Construction~A,
root-system, theta-series, and Sphere-Packing-Lean conventions.

%% -----------------------------------------------------------------------
\section{Mathematical Background}
\label{sec:background}

\subsection{The Extended Hamming Code}

Let $C$ be the extended binary Hamming code of length eight~\cite{Hamming1950}.
Classically, this is the unique Type~II binary code of length
eight~\cite{MacSloane1977,ConwaySloane1999}.  It has parameters $[8,4,4]$:
length eight, dimension four, and minimum distance four.  Its weight enumerator
is~\cite{MacSloane1977}
\[
  W_C(x,y) = x^8 + 14x^4y^4 + y^8.
\]
In Lean, $C$ is defined by an explicit
parity-check matrix over $\F_2 = \Z/2\Z$, and the code properties are proved
from that definition.  The formalization records the column ordering of the
matrix, since many classical parity-check matrices for the same code differ by
coordinate permutation.

The promoted code-theoretic Lean declarations are:
\begin{lstlisting}[style=lean]
CodeLatticeE8.Code.extendedHamming8_isLinearCode_4_4
CodeLatticeE8.Code.extendedHamming8_doublyEven_mem
CodeLatticeE8.Code.extendedHamming8_selfDual
CodeLatticeE8.Code.extendedHamming8_typeII
CodeLatticeE8.Code.extendedHamming8_unique_up_to_equivalence
\end{lstlisting}
The first packages the $[8,4,4]$ parameters, the second records the doubly-even
property used by Construction~A, the next two prove self-duality and Type~II
status, and the final theorem states that any binary linear $[8,4,4]$ code is
coordinate-permutation equivalent to the chosen extended Hamming
code~\cite{MacSloane1977,HuffmanPless2003}.

\subsection{Construction A}

For a binary code $C \subseteq (\Z/2\Z)^n$, Construction~A defines a subgroup
of integer vectors~\cite{LeechSloane1971,ConwaySloane1999}:
\[
  A(C) = \bigl\{\, z \in \Z^n \;\bigm|\; z \bmod 2 \in C \,\bigr\}.
\]
For the extended Hamming code, this produces an integer lattice in
$\Z^8$~\cite{LeechSloane1971}.  The unscaled integer model has minimal squared
norm~$4$.  The usual E8 normalization divides the quadratic form by~$2$,
equivalently scaling all coordinates by $1/\sqrt{2}$, so the standard E8
minimal squared norm is~$2$~\cite{ConwaySloane1999}.

The formalization uses the unscaled integer model for the direct Construction~A
definition and then records the scaled form explicitly.  This avoids hiding the
factor of two that appears in theta coefficients, shell counts, and Gram
determinants.

\subsection{E8 Metric and Root Data}

The standard E8 root system has $240$ roots, each of squared norm~$2$~\cite{ConwaySloane1999}.
In the Construction~A integer coordinates, these correspond to vectors of
unscaled squared norm~$4$.  The formalization proves both the minimum norm
statement and the exact cardinality of the short-vector shell.

It then relates the short-vector shell to an independent root-list model in
the repository.  This gives a second confirmation that the code-built lattice
has the expected E8 root structure, beyond its determinant and parity
properties.

\subsection{Theta Series and $E_4$}

The theta series of E8 is the generating function for shell
counts~\cite{ConwaySloane1999}:
\[
  \Theta_{E_8}(q) = \sum_{v \in E_8} q^{\|v\|^2/2}.
\]
For the standard E8 lattice, the classical formula is~\cite{ConwaySloane1999}:
\[
  \Theta_{E_8}(q) = E_4(q) = 1 + 240\sum_{n>0}\sigma_3(n)\,q^n.
\]
In the Construction~A integer coordinates, coefficient~$n$ corresponds to the
unscaled shell condition $\mathrm{sqNorm}(z) = 4n$.  The Lean proof keeps this
conversion explicit.  It first identifies the Construction~A coefficient with a
Hamming weight-enumerator convolution, then transports shell counts to the E8
theta modular form, and finally compares that modular form with~$E_4$.

%% -----------------------------------------------------------------------
\section{Formalization Architecture}
\label{sec:architecture}

The development is organized around a small number of interfaces.

\subsection{Code and Construction~A Layer}

The generic binary-code and Construction~A definitions live in:
\begin{lstlisting}[style=lean]
CodeLatticeE8/Code/Binary.lean
CodeLatticeE8/Code/Equivalence.lean
CodeLatticeE8/Code/LinearCode.lean
CodeLatticeE8/Code/Hamming844Basic.lean
CodeLatticeE8/Code/Hamming844WeightEnumerator.lean
CodeLatticeE8/Code/Hamming844.lean
CodeLatticeE8/Code/Hamming844Permutation.lean
CodeLatticeE8/Code/Hamming844Systematic.lean
CodeLatticeE8/Code/Hamming844Uniqueness.lean
CodeLatticeE8/ConstructionA/Basic.lean
CodeLatticeE8/ConstructionA/Norm.lean
CodeLatticeE8/ConstructionA/Even.lean
CodeLatticeE8/ConstructionA/TypeII.lean
\end{lstlisting}
These files define binary vectors, Hamming weight, reduction mod~$2$, the
extended Hamming code, weight distribution, the $[8,4,4]$ parameter package,
Construction~A as an integer subgroup, mod-four norm divisibility for doubly-even
codes, the general Type~II scaled evenness and self-duality theorem, and
uniqueness up to code equivalence.

\subsection{Lattice and Root Layer}

The Construction~A lattice, its explicit basis, and its E8 root data are
organized through:
\begin{lstlisting}[style=lean]
CodeLatticeE8/E8/HammingConstruction.lean
CodeLatticeE8/E8/Basis.lean
CodeLatticeE8/E8/Span.lean
CodeLatticeE8/E8/Gram.lean
CodeLatticeE8/E8/Determinant.lean
CodeLatticeE8/E8/ShortVectors.lean
CodeLatticeE8/E8/Roots.lean
CodeLatticeE8/E8/RootBridge.lean
CodeLatticeE8/E8/CartanBridge.lean
CodeLatticeE8/E8/WeylReflections.lean
\end{lstlisting}
The root-list, short-vector, root-bridge, Cartan-bridge, and basic Weyl
reflection layers are now part of the package.  Direct
Sphere-Packing-Lean comparisons remain optional and belong under the
\code{CodeLatticeE8SPL} root.  The theorem index is
\code{CodeLatticeE8/Publication/TheoremIndex.lean}.

\subsection{Theta-Series Layer}

The theta-series material is split into a standalone finite layer and an
analytic/Sphere-Packing-Lean layer.

The standalone package records the coefficient normalization and the finite
Hamming weight-contribution table:
\begin{lstlisting}[style=lean]
CodeLatticeE8/Theta/Sigma.lean
CodeLatticeE8/ConstructionA/ThetaLift.lean
CodeLatticeE8/ConstructionA/ThetaConvolution.lean
CodeLatticeE8/E8/ThetaCoefficients.lean
CodeLatticeE8/E8/ThetaTable.lean
CodeLatticeE8/E8/ThetaSeries.lean
\end{lstlisting}
The core finite table theorem is:
\begin{lstlisting}[style=lean]
CodeLatticeE8.E8.thetaTableCoeff_eq_e4Coeff_of_le_six
\end{lstlisting}
It states that a finite Hamming weight-contribution table agrees with the
normalized $E_4$ coefficients through $q^6$.  This is a standalone table
check, not a semantic shell-count theorem and not the full analytic identity.

The standalone package also contains the all-shell Construction~A convolution
theorem:
\begin{lstlisting}[style=lean]
CodeLatticeE8.E8.thetaShellCount_eq_convolution
\end{lstlisting}
It expresses the semantic Construction~A shell count as the Hamming
weight-distribution convolution for every unscaled squared norm.

The analytic layer has a root, \code{CodeLatticeE8SPL}, that
contains the completed all-coefficient bridge:
\begin{lstlisting}[style=lean]
CodeLatticeE8.SPL.thetaE8_MF_eq_E4
CodeLatticeE8.SPL.thetaE8_MF_qExpansion_coeff_eq_hammingThetaConvolutionCoeff
CodeLatticeE8.SPL.hammingThetaConvolutionCoeff_eq_e4Coeff
CodeLatticeE8.SPL.thetaSeries_hammingE8_eq_e4Series
\end{lstlisting}

\subsection{Sphere-Packing-Lean Bridge}

The bridge to Sphere-Packing-Lean uses explicit matrices and linear maps rather
than an abstract uniqueness theorem.  The local development reproduces the E8
basis matrix and compares Gram forms.  The imported bridge then gives a
coordinate map from the local Construction~A model to
$\code{Submodule.E8}\;R$.

This layer matters because Sphere-Packing-Lean's analytic result is stated for
its own E8 object.  The present formalization provides the code-to-lattice
identification needed to view the extended Hamming code as a finite entrance to
that endpoint.

%% -----------------------------------------------------------------------
\section{Main Formal Results}
\label{sec:results}

\subsection{The Hamming Code Is the Type~II Code of Length Eight}

The Lean theorem
\begin{lstlisting}[style=lean]
theorem extendedHamming8_typeII :
    IsTypeII extendedHamming8
\end{lstlisting}
states that the explicitly defined code is self-dual and doubly even.  The
parameter theorem is:
\begin{lstlisting}[style=lean]
theorem extendedHamming8_isLinearCode_4_4 :
    IsLinearCode extendedHamming8 4 4
\end{lstlisting}
and the promoted doubly-even theorem used directly by Construction~A is:
\begin{lstlisting}[style=lean]
theorem extendedHamming8_doublyEven_mem
    (v : BinaryVector 8) (hv : v in extendedHamming8) :
    4 divides hammingWeight v
\end{lstlisting}
This is the code-theoretic parity condition that makes the Construction~A norm
divisible by four.

The uniqueness theorem is:
\begin{lstlisting}[style=lean]
theorem extendedHamming8_unique_up_to_equivalence
    (C : BinaryLinearCode 8) (hC : IsLinearCode C 4 4) :
    CodeEquivalent C extendedHamming8
\end{lstlisting}
This says that the chosen parity-check representation is not an arbitrary
example: it represents, up to coordinate permutation, the unique binary linear
code with parameters $[8,4,4]$~\cite{MacSloane1977,ConwaySloane1999}.

\subsection{The Construction~A Lattice Has the E8 Gram Data}

The Construction~A subgroup comes with an explicit integer basis
\code{hammingConstructionBasis}.  The promoted spanning theorem is:
\begin{lstlisting}[style=lean]
theorem hammingConstructionASubmodule_eq_span :
    hammingConstructionASubmodule =
      Submodule.span Z (Set.range hammingConstructionBasis)
\end{lstlisting}
It states that the Construction~A lattice is exactly the $\Z$-span of this
basis.  The displayed Gram matrix is named \code{hammingConstructionGram}, and
its entrywise correctness theorem is:
\begin{lstlisting}[style=lean]
theorem hammingConstructionGram_eq :
    forall i j : Fin 8,
      hammingConstructionGram i j =
        sum k, hammingConstructionBasis i k * hammingConstructionBasis j k
\end{lstlisting}

Since the E8 normalization divides the quadratic form by~$2$, the scaled Gram
matrix is \code{hammingConstructionScaledGram : Matrix (Fin 8) (Fin 8) Q},
with diagonal theorem:
\begin{lstlisting}[style=lean]
theorem hammingConstructionScaledGram_diag (i : Fin 8) :
    hammingConstructionScaledGram i i = 2
\end{lstlisting}

The determinant computation is also in the package:
\begin{lstlisting}[style=lean]
theorem hammingConstructionGram_det :
    hammingConstructionGram.det = 256

theorem hammingConstructionScaledGram_det :
    hammingConstructionScaledGram.det = 1
\end{lstlisting}
These are proved through the explicit basis matrix rather than by expanding the
full $8\times 8$ determinant directly: the basis matrix is block upper
triangular with a $4\times 4$ determinant~$-1$ block and a $2I_4$ block, so
its determinant is~$-16$; the raw Gram determinant is its square, and the
scaled Gram determinant is obtained by \code{Matrix.det\_smul}.

Together with the divisibility theorem
\begin{lstlisting}[style=lean]
theorem hammingConstructionA_sqNorm_dvd_four
    (z : Fin 8 -> Z) (hz : z in hammingConstructionA) :
    (4 : Z) divides sqNorm z
\end{lstlisting}
these results give the integral basis, Gram, determinant, and evenness
package.

\subsection{Minimum Norm and Kissing Number}

The scaled minimum norm theorem is:
\begin{lstlisting}[style=lean]
theorem hammingConstructionA_scaledMinSqNorm :
    (forall z, z in hammingConstructionA -> z != 0 -> 2 <= scaledSqNorm z) and
    (exists z, z in hammingConstructionA and z != 0 and scaledSqNorm z = 2)
\end{lstlisting}

The exact short-vector count is:
\begin{lstlisting}[style=lean]
theorem hammingConstructionA_short_vector_count :
    Set.ncard hammingE8ShortShell = 240
\end{lstlisting}

Completeness is supplied by:
\begin{lstlisting}[style=lean]
theorem shortShell_mem_shortShellVectorList
    {z : Fin 8 -> Z}
    (hz : z in hammingE8ShortShell) :
    z in shortShellVectorList
\end{lstlisting}
Thus the Construction~A lattice has the E8 kissing number~\cite{ConwaySloane1999}.
In unscaled coordinates, the short shell is $\mathrm{sqNorm}(z) = 4$; in the
scaled E8 normalization, this is squared norm~$2$.  The finite cardinality
theorem is proved by private finite parametrizations of the weight-zero and
weight-four Hamming classes.

\subsection{The Root and Cartan Bridges}

The short-vector shell is compared to an independently defined E8 root list.
The permutation theorem is:
\begin{lstlisting}[style=lean]
theorem shortShell_perm_rootList :
    (shortShellVectorList.map shortVectorToRootCoords).Perm rootList
\end{lstlisting}
This theorem says that the Construction~A short vectors, after the explicit
doubled-coordinate bridge map, are precisely the root-list model.  The root
list itself is characterized by:
\begin{lstlisting}[style=lean]
theorem mem_rootList_iff_isE8Root (v : Fin 8 -> Z) :
    v in rootList iff IsE8Root v
\end{lstlisting}

The Cartan bridge is:
\begin{lstlisting}[style=lean]
theorem gramCartan_congruence :
    e8BasisChangeMatrix.transpose * hammingConstructionGram *
        e8BasisChangeMatrix =
      2 * e8CartanMatrix
\end{lstlisting}
It states that an explicit unimodular transition matrix carries the
Construction~A Gram matrix to the standard E8 Cartan form~\cite{Bourbaki2002}.
This provides a matrix-level identification of the code-built lattice with the
root-system presentation.  The congruence itself is proved by ordinary kernel
reduction over the 64 matrix entries.  The Cartan determinant is proved by
nested cofactor expansion, reducing the calculation to small $7\times 7$ kernel
\lean{decide} checks.

The package also records the first Weyl-theoretic closure property of
this root model.  In doubled coordinates, reflection through a root is~\cite{Bourbaki2002}:
\[
  \sigma_r(v) = v - \frac{\langle v,\, r\rangle}{4}\,r.
\]
The theorem
\begin{lstlisting}[style=lean]
theorem reflect_mem_rootList
    {r v : Fin 8 -> Z}
    (hr : r in rootList) (hv : v in rootList) :
    reflect r v in rootList
\end{lstlisting}
shows that the $240$-root list is closed under Weyl reflections.  The companion
theorems \lean{reflect\_reflect\_of\_mem} and \lean{reflect\_self\_eq\_neg}
record involutivity and reflection of a root through itself.  Closure and
involutivity are proved from the semantic root predicate \lean{IsE8Root}: the
key structural lemmas show that doubled inner products are divisible by~$4$,
that reflection preserves the semantic root conditions, and that reflection is
involutive.

There are two integer dot products in play.  The Cartan bridge uses
\lean{CodeLatticeE8.E8.intDot} in Construction~A coordinates.  The
Weyl-reflection module uses \lean{CodeLatticeE8.E8.Roots.dot} in doubled
root-list coordinates.  The Hadamard bridge relating these models scales squared
norms by~$8$, as recorded in
\lean{CodeLatticeE8.E8.RootBridge.hadamard\_sqNorm\_scale}; this is why a
Construction~A short vector with $\mathrm{sqNorm} = 4$ maps to a doubled root
with $\mathrm{Roots.normSq} = 8$, while both represent an E8 vector of squared
norm~$2$.

\subsection{Theta Coefficients and the $\Theta_{E_8} = E_4$ Bridge}

The standalone package currently exposes the theta comparison in two finite
layers.  First, it proves the semantic shell statements:
\begin{lstlisting}[style=lean]
theorem semanticThetaCoeff_eq_e4Coeff_zero :
    Set.ncard (hammingE8Shell 0) = e4Coeff 0

theorem semanticThetaCoeff_eq_e4Coeff_one :
    Set.ncard (hammingE8Shell 4) = e4Coeff 1
\end{lstlisting}
The $q^1$ result is exactly the kissing-number theorem restated as the first
nonconstant theta coefficient.

Second, the core package records a finite Hamming weight-contribution table and
proves:
\begin{lstlisting}[style=lean]
theorem thetaTableCoeff_eq_e4Coeff_of_le_six
    (n : Nat) (hn : n <= 6) :
    hammingThetaTableCoeff n = e4Coeff n
\end{lstlisting}
This is a standalone finite table theorem.  The semantic standalone statement
is the all-shell Construction~A convolution theorem
\lean{CodeLatticeE8.E8.thetaShellCount\_eq\_convolution}; the SPL-facing
package then identifies that semantic coefficient function with $E_4$ for
every~$n$.

The SPL-facing theorem package uses a semantic coefficient function
$\code{hammingThetaConvolutionCoeff}(n)$ for the Construction~A coefficient at
theta index~$n$.  The completed all-coefficient theorem is:
\begin{lstlisting}[style=lean]
theorem hammingThetaConvolutionCoeff_eq_e4Coeff (n : Nat) :
    hammingThetaConvolutionCoeff n = Theta.e4Coeff n
\end{lstlisting}
The formal power-series target is:
\begin{lstlisting}[style=lean]
theorem CodeLatticeE8.SPL.thetaSeries_hammingE8_eq_e4Series :
    thetaSeries = e4Series
\end{lstlisting}
The proof route factors through the modular-form theorem:
\begin{lstlisting}[style=lean]
theorem CodeLatticeE8.SPL.thetaE8_MF_eq_E4 :
    thetaE8_MF = E4
\end{lstlisting}
and the $q$-expansion bridge:
\begin{lstlisting}[style=lean]
theorem CodeLatticeE8.SPL.
  thetaE8_MF_qExpansion_coeff_eq_hammingThetaConvolutionCoeff
    (n : Nat) :
    (ModularFormClass.qExpansion (1 : R) thetaE8_MF).coeff n =
      (hammingThetaConvolutionCoeff n : C)
\end{lstlisting}
The mathematical content of the bridge is that the Fourier coefficient of the
rank-eight theta modular form is the cardinality of the E8 shell of squared
norm~$2n$, and that this shell is bijective with the Construction~A shell
$\mathrm{sqNorm}(z) = 4n$.

\subsection{Bridge to Sphere-Packing-Lean}

Sphere-Packing-Lean uses its own E8 basis and lattice object.  The local
formalization compares the Construction~A basis to that object by explicit
linear algebra.

The key point is not merely that both lattices have the same determinant or the
same number of roots.  The bridge supplies concrete maps between coordinate
models and proves the corresponding inner-product formula.  Thus the Hamming
Construction~A lattice is connected directly to the E8 object used by the
formalized analytic sphere-packing proof.

%% -----------------------------------------------------------------------
\section{Proof Ideas}
\label{sec:proofs}

\subsection{From the Code to the Lattice}

The Construction~A proof is deliberately concrete.  Membership in the lattice
is defined by reduction mod~$2$ into the Hamming code.  The explicit basis is
then shown to span exactly this subgroup.  The Gram matrix is computed from the
basis, and the determinant and divisibility properties are proved from the
matrix and code equations.

The package also contains the corresponding general Type~II
Construction~A theorem.  In Lean, the integer-coordinate theorem is
\lean{CodeLatticeE8.ConstructionA.typeII\_integer\_package}; the real scaled
dual statement is
\lean{CodeLatticeE8.ConstructionA.scaledRealDual\_eq\_self\_of\_selfDual}.
Thus the formal library now records the standard abstraction that Type~II
codes give even self-dual, hence unimodular, lattices after the conventional
$1/\sqrt{2}$ scaling~\cite{MacSloane1977,ConwaySloane1999}.  The paper still
keeps the rank-eight basis, determinant, and Cartan comparisons explicit,
because those are the auditable coordinate bridges needed to identify the Hamming
Construction~A model with the E8 lattice used elsewhere in the formalization.

\subsection{The Short-Vector Shell}

The short vectors split into the two familiar E8 families.  In Construction~A
integer coordinates, vectors of squared norm~$4$ arise from:
\begin{enumerate}
\item coordinate spikes $(\pm 2)$ in one coordinate, reducing to the zero
      codeword; and
\item sign choices on four coordinates supported on a weight-four Hamming
      codeword.
\end{enumerate}
The weight distribution $(1,14,1)$~\cite{MacSloane1977} explains the cardinality:
\[
  16 + 14 \times 16 = 240.
\]
The Lean proof packages this as an explicit finite list and proves both
soundness and completeness for the short-vector predicate.

\subsection{The Theta-Series Proof Plan}

The coefficient proof follows the classical theta-series route, but each
normalization step is explicit.

The Construction~A side first proves a convolution theorem.  Fixing a binary
residue word, the number of integer lifts with a given squared norm is a
product of one-dimensional even or odd lift counts.  Summing over Hamming
codewords and grouping by weight gives the function
$\code{hammingThetaConvolutionCoeff}$.

The standalone package proves the semantic $q^0$ and $q^1$ coefficients, the
finite Hamming contribution table through $q^6$, and the all-shell
Construction~A convolution theorem.  The SPL-facing layer completes
the all-coefficient comparison with~$E_4$: it constructs a rank-eight theta
modular form, proves its equality with~$E_4$, extracts $q$-expansion
coefficients as shell counts, and transports those shell counts to the Hamming
Construction~A model.  The final bridge proves that the analytic E8 shell and
the Construction~A shell are in bijection with the correct factor of two in the
norm:
\[
  \mathrm{sqNorm}_{\mathrm{CA}}(z) = 4n \iff \|\phi(z)\|^2 = 2n.
\]
Combining these facts yields the all-coefficient formula $1$ at $n=0$ and
$240\,\sigma_3(n)$ for $n>0$, recorded in
\lean{CodeLatticeE8.SPL.hammingThetaConvolutionCoeff\_eq\_e4Coeff}.

\subsection{Why an Explicit SPL Bridge Matters}

The theorem $\Theta_{E_8} = E_4$ identifies the theta series
abstractly, but the sphere-packing endpoint is phrased for the E8 model used
inside Sphere-Packing-Lean.  The explicit matrix and lattice bridges ensure
that the code-built lattice and the SPL lattice are not merely isomorphic in
principle; they are connected by named Lean maps and inner-product statements.

This is the central formalization lesson of the paper.  In ordinary
mathematical prose, one often moves freely among equivalent E8 presentations.
In Lean, every such move becomes a theorem.  The benefit is that the final
statement carries no hidden convention shift.

%% -----------------------------------------------------------------------
\section{Verification and Trust Boundary}
\label{sec:trust}

The repository is pinned to Lean~4.28.0.  The main theorem statements cited in
this paper should be checked by building the relevant Lean modules under the
pinned toolchain.

Some files use local \lean{maxHeartbeats} increases and style-linter
suppressions for large finite matrix reductions or long analytic tactic scripts.
These are engineering choices about proof-search budget and warning noise; the
proofs remain ordinary kernel-checked Lean terms and the trust report records
the absence of \lean{native\_decide} and \lean{Lean.trustCompiler} from the
publication spine.

Minimal checks for the local theorem index:
\begin{lstlisting}[style=shell]
lake build CodeLatticeE8
\end{lstlisting}
The standalone Lean modules imported by \code{CodeLatticeE8} do not import
Sphere-Packing-Lean.  Reviewers who want to check only this core layer without
resolving the optional SPL dependency should use the standalone Lake wrapper:
\begin{lstlisting}[style=shell]
cd CodeLatticeE8Standalone
lake build CodeLatticeE8
\end{lstlisting}
Checks for the optional SPL-facing theta-series route:
\begin{lstlisting}[style=shell]
lake build CodeLatticeE8SPL
\end{lstlisting}
The theorem names and files used in the paper are listed in
Appendix~\ref{app:theorems}.  Appendix~\ref{app:verification} records the
intended final verification checklist, including axiom reports for the theorem
declarations that form the paper spine.

%% -----------------------------------------------------------------------
\section{Related Work}
\label{sec:related}

The classical mathematical background is the code-lattice theory developed in
MacWilliams--Sloane~\cite{MacSloane1977}, Conway--Sloane~\cite{ConwaySloane1999},
and later lattice and coding-theory texts~\cite{HuffmanPless2003}.  The extended
Hamming code is the smallest Type~II binary code, and Construction~A applied to
it is the standard coding-theoretic construction of
E8~\cite{LeechSloane1971,ConwaySloane1999}.

Sphere-Packing-Lean~\cite{SPL2025} formalizes the analytic endpoint of the
dimension-eight sphere-packing theorem, following the Cohn--Elkies
linear-programming method~\cite{CohnElkies2003} and Viazovska's
construction~\cite{Viazovska2017}.  The present work is not a replacement for
that analytic formalization.  Instead, it supplies a formal combinatorial bridge
to the E8 object used there.

Recent Lean formalizations of coding theory, self-dual codes, and related
algebraic constructions provide adjacent infrastructure.  This paper differs in
focus: it follows one specific code all the way to the exceptional lattice and
then connects the result to root-system, theta-series, and sphere-packing
models.

%% -----------------------------------------------------------------------
\section{Conclusion}
\label{sec:conclusion}

The extended Hamming code is a compact finite object.  In this formalization it
serves as a front door to E8: from a parity-check matrix over $\F_2$, Lean
constructs a lattice, computes its Gram form, proves its minimal-vector
structure, identifies its root data, proves its theta-series formula, and
bridges it to the E8 lattice used by Sphere-Packing-Lean.

The result is a verified path through several classical E8 presentations.  It
also provides reusable infrastructure for future work: Golay-to-Leech
Construction~A, theta-series identities for other lattices, and physics-facing
code-lattice constructions such as $E_8\times E_8$ models.

%% -----------------------------------------------------------------------
%% Appendices
%% -----------------------------------------------------------------------
\appendix

\section{Conventions and Normalizations}
\label{app:conventions}

The paper uses the following models.

\begin{table}[h]
\centering
\renewcommand{\arraystretch}{1.35}
\begin{tabular}{@{}llll@{}}
\toprule
\textbf{Model} & \textbf{Ambient type} & \textbf{Short sq.\ norm} & \textbf{Role} \\
\midrule
Construction~A integer model &
  $\Z^8$ (via $\code{Fin}\ 8 \to \Z$) &
  $4$ &
  Direct output of Construction~A \\
Scaled E8 model &
  same basis, form $/2$ &
  $2$ &
  Standard E8 normalization \\
Doubled root model &
  $\Z^8$ &
  $8$ (doubled coords) &
  Root-list bridge \\
SPL model &
  $\code{EuclideanSpace}\ \R\ (\code{Fin}\ 8)$ &
  $2$ &
  Sphere-Packing-Lean convention \\
\bottomrule
\end{tabular}
\caption{Coordinate models and their normalization conventions.}
\label{tab:models}
\end{table}

The theta coefficient indexed by~$n$ counts the unscaled Construction~A shell
\[
  \mathrm{sqNorm}(z) = 4n.
\]
Under the bridge to the standard E8 normalization, this becomes
$\|\phi(z)\|^2 = 2n$.
Equivalently, the theta bridge uses the normalization
\[
  \mathrm{sqNorm}(z) = 4n \iff \|\phi(z)\|^2 = 2n.
\]
The Lean declaration \lean{CodeLatticeE8.SPL.sqNorm\_eq\_two\_mul\_norm\_sq}
records the underlying equality $\mathrm{sqNorm}(z) = 2\|\phi(z)\|^2$ for the
explicit bridge map.  This is the most important normalization in the
theta-series proof.

\section{Lean Theorem Index}
\label{app:theorems}

The following table gives the primary theorem names for the paper.
Displayed theorem statements in the body are lightly normalized for
readability.  The declaration names and file paths in the table are the
authoritative handles for checking the exact Lean statements.

\begin{longtable}{@{}p{5.3cm}p{5.4cm}p{3.5cm}@{}}
\toprule
\textbf{Mathematical claim} &
\textbf{Lean declaration} &
\textbf{File} \\
\midrule
\endhead
\bottomrule
\endfoot
Extended Hamming code has parameters $[8,4,4]$ &
  \lean{...extendedHamming8\_isLinearCode\_4\_4} &
  \lean{Code/Hamming844.lean} \\[2pt]
Extended Hamming code is doubly even &
  \lean{...extendedHamming8\_doublyEven\_mem} &
  \lean{Code/Hamming844.lean} \\[2pt]
Extended Hamming code is self-dual &
  \lean{...extendedHamming8\_selfDual} &
  \lean{Code/Dual.lean} \\[2pt]
Extended Hamming code is Type~II &
  \lean{...extendedHamming8\_typeII} &
  \lean{Code/Dual.lean} \\[2pt]
$[8,4,4]$ uniqueness &
  \lean{...extendedHamming8\_unique\_up\_to\_equivalence} &
  \lean{Code/Hamming844Uniqueness.lean} \\[2pt]
Construction~A minimum unscaled norm $4$ &
  \lean{...hammingConstructionA\_minSqNorm} &
  \lean{E8/HammingConstruction.lean} \\[2pt]
Construction~A norm divisibility &
  \lean{...hammingConstructionA\_sqNorm\_dvd\_four} &
  \lean{E8/HammingConstruction.lean} \\[2pt]
Generic Type~II integer package &
  \lean{...ConstructionA.typeII\_integer\_package} &
  \lean{ConstructionA/TypeII.lean} \\[2pt]
Generic Type~II real self-duality &
  \lean{...scaledRealDual\_eq\_self\_of\_selfDual} &
  \lean{ConstructionA/TypeII.lean} \\[2pt]
Explicit basis membership &
  \lean{...hammingConstructionBasis\_mem} &
  \lean{E8/Basis.lean} \\[2pt]
Explicit basis spans the lattice &
  \lean{...hammingConstructionASubmodule\_eq\_span} &
  \lean{E8/Span.lean} \\[2pt]
Unscaled Gram matrix entries &
  \lean{...hammingConstructionGram\_eq} &
  \lean{E8/Basis.lean} \\[2pt]
Scaled Gram diagonal has root norm $2$ &
  \lean{...hammingConstructionScaledGram\_diag} &
  \lean{E8/Gram.lean} \\[2pt]
Scaled Gram entries are integral &
  \lean{...hammingConstructionScaledGram\_int\_valued} &
  \lean{E8/Gram.lean} \\[2pt]
Minimum squared norm $2$ &
  \lean{...hammingConstructionA\_scaledMinSqNorm} &
  \lean{E8/Minimum.lean} \\[2pt]
Unscaled Gram determinant $= 256$ &
  \lean{...hammingConstructionGram\_det} &
  \lean{E8/Determinant.lean} \\[2pt]
Scaled Gram determinant $= 1$ &
  \lean{...hammingConstructionScaledGram\_det} &
  \lean{E8/Determinant.lean} \\[2pt]
Kissing number $240$ &
  \lean{...hammingConstructionA\_short\_vector\_count} &
  \lean{E8/ShortVectors.lean} \\[2pt]
Short-vector completeness &
  \lean{...RootBridge.shortShell\_mem\_shortShellVectorList} &
  \lean{E8/RootBridge.lean} \\[2pt]
E8 root-list characterization &
  \lean{...Roots.mem\_rootList\_iff\_isE8Root} &
  \lean{E8/Roots.lean} \\[2pt]
Root-list bridge &
  \lean{...RootBridge.shortShell\_perm\_rootList} &
  \lean{E8/RootBridge.lean} \\[2pt]
Cartan bridge &
  \lean{...gramCartan\_congruence} &
  \lean{E8/CartanBridge.lean} \\[2pt]
Weyl reflection closure &
  \lean{...WeylReflections.reflect\_mem\_rootList} &
  \lean{E8/WeylReflections.lean} \\[2pt]
$E_4$ coefficient normalization &
  \lean{CodeLatticeE8.Theta.e4Coeff} &
  \lean{Theta/Sigma.lean} \\[2pt]
Semantic theta $q^0$ shell &
  \lean{...semanticThetaCoeff\_eq\_e4Coeff\_zero} &
  \lean{E8/ThetaCoefficients.lean} \\[2pt]
Semantic theta $q^1$ shell &
  \lean{...semanticThetaCoeff\_eq\_e4Coeff\_one} &
  \lean{E8/ThetaCoefficients.lean} \\[2pt]
Hamming contribution table through $q^6$ &
  \lean{...thetaTableCoeff\_eq\_e4Coeff\_of\_le\_six} &
  \lean{E8/ThetaSeries.lean} \\[2pt]
All-shell Construction~A convolution &
  \lean{...thetaShellCount\_eq\_convolution} &
  \lean{E8/ThetaSeries.lean} \\[2pt]
E8 theta modular form $= E_4$ &
  \lean{CodeLatticeE8.SPL.thetaE8\_MF\_eq\_E4} &
  \lean{SPL/E8ThetaModular.lean} \\[2pt]
$q$-expansion to Hamming coefficient &
  \lean{...thetaE8\_MF\_qExpansion\_coeff\_eq\_...} &
  \lean{SPL/CoefficientBridge.lean} \\[2pt]
E8 representation formula &
  \lean{...hammingThetaConvolutionCoeff\_eq\_e4Coeff} &
  \lean{SPL/CoefficientBridge.lean} \\[2pt]
Formal theta series $=$ formal $E_4$ &
  \lean{...thetaSeries\_hammingE8\_eq\_e4Series} &
  \lean{SPL/CoefficientBridge.lean} \\
\end{longtable}

\noindent
All \lean{CodeLatticeE8.*} declarations use the full namespace prefix.
The optional SPL root \code{CodeLatticeE8SPL.lean} imports
\code{CodeLatticeE8.SPL.TheoremIndex}, which lists the completed analytic,
$q$-expansion, shell-transport, and formal power-series theorem chain.

The standalone package currently has no live \lean{native\_decide} dependency
in the paper spine.

\section{Hamming Code Details}
\label{app:hamming}

The extended Hamming code is defined by an explicit parity-check matrix over
$\F_2$~\cite{Hamming1950,MacSloane1977}.  Four parity equations determine a
four-dimensional subspace of $\F_2^8$.  The formal development proves:
\begin{enumerate}
\item the code has $16$ codewords;
\item the minimum nonzero Hamming weight is $4$;
\item every codeword has weight divisible by $4$;
\item the code is self-dual;
\item the weight distribution is $(1,14,1)$;
\item every binary linear $[8,4,4]$ code is equivalent to it~\cite{MacSloane1977}.
\end{enumerate}

The weight enumerator is~\cite{MacSloane1977}:
\[
  W_C(x,y) = x^8 + 14x^4y^4 + y^8.
\]
This polynomial is the finite combinatorial input for the Construction~A
theta-series computation.

\section{Construction~A Basis and Gram Data}
\label{app:gram}

The Construction~A lattice is~\cite{LeechSloane1971,ConwaySloane1999}:
\[
  A(C) = \bigl\{\, z \in \Z^8 \;\bigm|\; z \bmod 2 \in C_{\mathrm{Hamming}} \,\bigr\}.
\]
The clean formalization gives an explicit basis \lean{hammingConstructionBasis}.
The theorem \lean{hammingConstructionASubmodule\_eq\_span} proves that this
basis spans the Construction~A submodule.  From this basis it computes the Gram
matrix \lean{hammingConstructionGram} and proves:
\[
  G_{ij} = \sum_{k} b_i(k)\,b_j(k),
\]
where $b_i$ denotes the $i$-th basis vector.

Because the E8 normalization divides the quadratic form by~$2$, the clean
package defines the rational scaled Gram matrix
\lean{hammingConstructionScaledGram} and proves $\tilde{G}_{ii} = 2$ for every
diagonal entry.

The determinant proof is cleanly packaged.  The basis matrix determinant is
proved by a block-upper-triangular decomposition:
\[
  \det(B) = -16, \quad \det(G) = (-16)^2 = 256, \quad
  \det(\tilde{G}) = 2^{-8}\cdot 256 = 1.
\]
The same basis supports the Cartan comparison.  An explicit unimodular
transition matrix~$P$ (with $\det P = \pm 1$) carries the Construction~A Gram
form to the standard E8 Cartan form~\cite{Bourbaki2002,ConwaySloane1999}:
\[
  P^\top G P = 2\,\mathrm{Cartan}_{E_8}.
\]

\section{Short Vectors and Roots}
\label{app:roots}

In unscaled Construction~A coordinates, the short shell is:
\[
  \bigl\{\, z \in A(C) \;\bigm|\; \mathrm{sqNorm}(z) = 4 \,\bigr\}.
\]
The formal proof classifies this shell by Hamming residue.

\medskip
\noindent\textit{Zero residue (type-1 vectors).}
Each coordinate spike $(\ldots,\pm 2,\ldots)$ contributes $8\times 2 = 16$
vectors.

\medskip
\noindent\textit{Weight-four residue (type-2 vectors).}
Each weight-four codeword contributes sign choices $(\pm 1,\pm 1,\pm 1,\pm 1)$
on its support, giving $2^4 = 16$ choices.  Since there are $14$ weight-four
codewords, this contributes $14\times 16 = 224$ vectors.

\medskip
\noindent The total is $16 + 224 = 240$~\cite{ConwaySloane1999}.

The Lean development proves not only this count, but also that the explicit
short-vector list is complete and maps to the repository's E8 root list.

\section{Theta-Series Details}
\label{app:theta}

The Construction~A theta coefficient at index~$n$ is the cardinality
of~\cite{ConwaySloane1999,LeechSloane1971}:
\[
  \bigl\{\, z \in A(C) \;\bigm|\; \mathrm{sqNorm}(z) = 4n \,\bigr\}.
\]

The package exposes two finite theta layers.

\paragraph{Semantic shells.}
\lean{CodeLatticeE8/E8/ThetaCoefficients.lean} records:
\begin{lstlisting}[style=lean]
CodeLatticeE8.E8.semanticThetaCoeff_eq_e4Coeff_zero
CodeLatticeE8.E8.semanticThetaCoeff_eq_e4Coeff_one
\end{lstlisting}
The $q^1$ theorem reuses the proved short-vector count.

\paragraph{Finite Hamming weight-contribution table.}
\lean{CodeLatticeE8/E8/ThetaTable.lean} records the table through $q^6$.  For
a codeword of weight~$w$, odd coordinates contribute square $1$ or~$9$, and
even coordinates contribute square $0$, $4$, or~$16$ in the verified range
$\mathrm{sqNorm} \le 24$.  The weight multiplicities are~\cite{MacSloane1977}:
\[
  A_0 = 1, \quad A_4 = 14, \quad A_8 = 1.
\]
The resulting table values are checked against the displayed $E_4$ coefficients
in \lean{thetaTableCoeff\_eq\_e4Coeff\_of\_le\_six}.

The structural Construction~A convolution theorem
\lean{thetaShellCount\_eq\_convolution} gives the semantic shell count as the
Hamming weight-distribution convolution for every unscaled squared norm.

\paragraph{Optional analytic layer.}
The optional SPL-facing analytic comparison uses the E8 theta modular form
\lean{thetaE8\_MF}.  Its $q$-expansion coefficient at~$n$ is identified with
the standard E8 shell $\|v\|^2 = 2n$.  The modular-form theorem proves
$\lean{thetaE8\_MF} = E_4$, giving the coefficient formula:
\[
  \text{coeff}_0 = 1, \qquad
  \text{coeff}_n = 240\,\sigma_3(n) \quad \text{for } n > 0.
\]

\section{Sphere-Packing-Lean Bridge}
\label{app:spl}

The local code-built lattice and the Sphere-Packing-Lean E8 lattice use
different coordinate models.  The formal bridge records:
\begin{enumerate}
\item a local reproduction of the relevant SPL-shaped E8 basis matrix;
\item Gram matrix comparison with the local Construction~A scaled Gram form;
\item a coordinate transition map;
\item a lattice equivalence with $\code{Submodule.E8}\ R$;
\item an inner-product preservation formula.
\end{enumerate}
This bridge is what allows the manuscript to say that the extended Hamming code
constructs the E8 lattice used at the formalized sphere-packing endpoint.

\section{Verification and Reproducibility}
\label{app:verification}

The Lean artifacts are checked in three layers, matching the dependency
boundaries used in the paper.

\paragraph{Standalone Hamming-to-E8 package.}
The main package is checked by:
\begin{lstlisting}[style=shell]
lake build CodeLatticeE8
\end{lstlisting}
For a mathlib-only checkout, the standalone wrapper gives the same root without
the Sphere-Packing-Lean dependency:
\begin{lstlisting}[style=shell]
cd CodeLatticeE8Standalone
lake build CodeLatticeE8
cd ..
\end{lstlisting}

\paragraph{Sphere-Packing-Lean-facing bridge.}
The analytic theta identity and the comparison with the E8 object used in
Sphere-Packing-Lean are checked by:
\begin{lstlisting}[style=shell]
lake build CodeLatticeE8SPL
\end{lstlisting}

\paragraph{Whole repository.}
The broader repository check is:
\begin{lstlisting}[style=shell]
lake build
\end{lstlisting}

The paper cites theorem declarations rather than informal file names whenever a
precise formal claim is involved.  Appendix~\ref{app:theorems} gives the
declaration/file map.  The trust report
\lean{Sources/CodeLatticeE8\_Trust\_Report.md} records the corresponding axiom
audit.  The expected axiom profile for the main paper-spine theorems is the
standard Lean/Mathlib set
\[
  \{\texttt{propext},\; \texttt{Classical.choice},\; \texttt{Quot.sound}\}.
\]
Some lower-level finite reductions may report \lean{Lean.ofReduceBool}; this is
Lean's kernel-checked boolean-reduction certificate, not a compiler-trust
assumption.  The standalone root is additionally checked to contain no
\lean{sorry}, \lean{admit}, fake axioms, \lean{unsafe def},
\lean{native\_decide}, or \lean{Lean.trustCompiler}.

\section{Bibliographic Roles}
\label{app:bibcheck}

The cited sources play the following roles in the paper.

\begin{itemize}
\item Hamming's original paper~\cite{Hamming1950} supplies the historical
      coding-theory origin of the Hamming code.
\item MacWilliams--Sloane~\cite{MacSloane1977} and
      Huffman--Pless~\cite{HuffmanPless2003} are the coding-theory references
      for the extended Hamming code, weight enumerators, self-duality, Type~II
      codes, and uniqueness up to coordinate permutation.
\item Leech--Sloane~\cite{LeechSloane1971} and
      Conway--Sloane~\cite{ConwaySloane1999} are the lattice references for
      Construction~A, the Hamming-code construction of E8, the E8 theta series,
      and the standard normalizations of even unimodular lattices.
\item Bourbaki~\cite{Bourbaki2002} is used for the E8 root-system and Cartan
      matrix conventions.
\item Cohn--Elkies~\cite{CohnElkies2003}, Viazovska~\cite{Viazovska2017}, and
      the Sphere-Packing-Lean report~\cite{SPL2025} provide the analytic
      sphere-packing endpoint and its Lean formalization.
\item The Error Correction Zoo entry~\cite{ECZ} is cited as a secondary
      orientation source for the named $[8,4,4]$ extended Hamming code.
\end{itemize}

%% -----------------------------------------------------------------------
\bibliographystyle{plain}
\begin{thebibliography}{99}

\bibitem{Hamming1950}
R.~W. Hamming.
\newblock Error detecting and error correcting codes.
\newblock \textit{Bell System Technical Journal}, 29(2):147--160, 1950.

\bibitem{LeechSloane1971}
J.~Leech and N.~J.~A. Sloane.
\newblock Sphere packings and error-correcting codes.
\newblock \textit{Canadian Journal of Mathematics}, 23(4):718--745, 1971.

\bibitem{MacSloane1977}
F.~J. MacWilliams and N.~J.~A. Sloane.
\newblock \textit{The Theory of Error-Correcting Codes}.
\newblock North-Holland Mathematical Library. North-Holland, Amsterdam, 1977.

\bibitem{HuffmanPless2003}
W.~C. Huffman and V.~Pless.
\newblock \textit{Fundamentals of Error-Correcting Codes}.
\newblock Cambridge University Press, Cambridge, 2003.

\bibitem{ConwaySloane1999}
J.~H. Conway and N.~J.~A. Sloane.
\newblock \textit{Sphere Packings, Lattices and Groups}.
\newblock Grundlehren der Mathematischen Wissenschaften 290. Springer-Verlag,
  New York, 3rd edition, 1999.

\bibitem{Bourbaki2002}
N.~Bourbaki.
\newblock \textit{Lie Groups and Lie Algebras}, Chapters~4--6.
\newblock Springer-Verlag, Berlin, 2002.

\bibitem{CohnElkies2003}
H.~Cohn and N.~Elkies.
\newblock New upper bounds on sphere packings~{I}.
\newblock \textit{Annals of Mathematics}, 157(2):689--714, 2003.

\bibitem{Viazovska2017}
M.~Viazovska.
\newblock The sphere packing problem in dimension~8.
\newblock \textit{Annals of Mathematics}, 185(3):991--1015, 2017.

\bibitem{SPL2025}
S.~Hariharan, C.~Birkbeck, S.~Lee, H.~K.~G. Ma, B.~Mehta, A.~Poiroux,
  and M.~Viazovska.
\newblock A milestone in formalization: The sphere packing problem in
  dimension~8.
\newblock \textit{arXiv:2604.23468}, 2025.

\bibitem{ECZ}
Error Correction Zoo.
\newblock $[8,4,4]$ extended {H}amming code.
\newblock \url{https://errorcorrectionzoo.org/c/hamming844}.

\end{thebibliography}

\end{document}
